\referenceTitle % Команда выведет РЕФЕРАТ жирным по центру и скроет номер страницы

\MakeUppercase{\taskNameFull} : дипломная работа / \studentShort. -- Минск : БГУИР, \targetYear. -- п.з. -- \getpagerefnumber{LastPage}~с., чертежей (плакатов) -- 6 л. формата А1.

\vspace{1em}

КЛЮЧЕВЫЕ СЛОВА: СЖАТИЕ ИЗОБРАЖЕНИЙ, ВЕЙВЛЕТ-ПРЕОБРАЗОВАНИЕ, СВЕРТОЧНЫЕ НЕЙРОННЫЕ СЕТИ, СВЕРХРАЗРЕШЕНИЕ, АСИММЕТРИЧНЫЕ ВЫЧИСЛЕНИЯ, EDGE-TO-CLOUD.

\vspace{1em}

Предмет исследования -- асимметричный гибридный алгоритм сжатия и восстановления растровых изображений.

Цель работы -- разработка, программная реализация и исследование гибридного программного комплекса, использующего вейвлет-декомпозицию для снижения битрейта и сверточные нейронные сети для восстановления разрешения.

В дипломной работе проведен анализ ограничений классических методов сжатия и изучен математический аппарат вейвлет-преобразования. Разработана и обоснована структурная схема асимметричной архитектуры Edge-to-Cloud, переносящей вычислительную нагрузку с мобильного устройства на сервер. Реализован математический кодер с экстракцией гибридных метаданных и интеллектуальный декодер с тайловой маршрутизацией памяти на базе моделей ESPCN и FSRCNN. Приводится расчет экономической эффективности разработки программного комплекса.

В результате экспериментов доказано, что разработанный алгоритм обеспечивает экстремальное снижение битрейта (до 0,1--0,2~бит/пиксель) при сохранении высоких показателей структурного сходства (SSIM 0,85--0,92), превосходя классические стандарты сжатия графической информации.